\section{Abstract}

Este estudio presenta un análisis comparativo de redes neuronales (\textbf{ResNet50 y VGG-16}) para la clasificación biométrica de perfiles de riesgo, eliminando el sesgo de entorno (\textit{shortcut learning}). Mediante \textbf{MTCNN}, se aislaron los óvalos faciales para procesar rasgos morfológicos y no fondos institucionales. Para mitigar sesgos algorítmicos, se implementó un \textbf{balanceo demográfico (4-way sampling)} usando \textbf{DeepFace} para generar metadatos étnicos, neutralizando prejuicios estadísticos entre sujetos blancos y negros de los datasets \textbf{Illinois DOC y LFW}. Los modelos alcanzaron una precisión del \textbf{97\%}, y la validación con \textbf{Grad-CAM} confirmó que la atención se focaliza en la estructura ósea y orbital. Finalmente, la auditoría con \textbf{sujetos externos} demostró que el equilibrio poblacional y la limpieza visual son imperativos para una clasificación justa y técnicamente robusta, capaz de generalizar más allá del entrenamiento.


\textbf{Palabras clave:} Biometría Forense, Deep Learning, ResNet50, VGG-16, Grad-CAM, Mitigación de Sesgo, MTCNN.

\section{Introducción}

La posibilidad de vincular la fisonomía humana con la conducta es una interrogante que ha persistido desde las teorías del atavismo de \textbf{Cesare Lombroso} hasta los estudios estadísticos de \textbf{Charles Goring}. En la actualidad, este debate cobra una relevancia renovada gracias a las \textbf{Redes neuronales}, que permite sustituir las observaciones subjetivas por un análisis matemático de patrones biométricos. El aporte fundamental de este trabajo consiste en proponer un marco técnico diseñado para evaluar si los sistemas de visión artificial pueden identificar rasgos faciales de manera objetiva, eliminando los prejuicios humanos que históricamente invalidaron este tipo de investigaciones.

El desafío teórico de este dominio reside en la intersección entre la biometría predictiva y la ética algorítmica, enfrentándose al fenómeno crítico del \textit{shortcut learning} (aprendizaje de atajos). Históricamente, los modelos han tendido a aprender correlaciones espurias, asociando la criminalidad a elementos del entorno como los fondos carcelarios o la iluminación de las fichas policiales. Por ello, la complejidad del problema no radica solo en la clasificación, sino en lograr que una red neuronal profunda aísle la morfología facial del "ruido" ambiental y demográfico, un reto de ingeniería de datos que busca garantizar que el aprendizaje sea puramente fisonómico.

Desde una perspectiva técnica, este proyecto se define como un problema de \textbf{aprendizaje supervisado} orientado a la clasificación binaria. A través de arquitecturas de vanguardia como \textbf{ResNet50} y \textbf{VGG-16}, el sistema categoriza a los sujetos en las clases de \textbf{"Riesgo"} y \textbf{"No Riesgo"} mediante la extracción de vectores de características de alta dimensionalidad. Más allá de la precisión predictiva, el enfoque se centra en la explicabilidad, utilizando mapas de activación para entender qué regiones faciales ponderan los modelos en su toma de decisiones, validando así la solidez científica del proceso.

A lo largo de las siguientes secciones, el lector encontrará un análisis detallado que comienza con el estudio exploratorio de datos (\textbf{EDA}) para auditar sesgos iniciales en los conjuntos de datos. Posteriormente, se describe el proceso de preprocesamiento facial mediante \textbf{MTCNN} y la configuración específica de las redes neuronales empleadas. El artículo concluye con una auditoría exhaustiva de los resultados mediante \textbf{matrices de confusión} y una fase de explicabilidad visual con \textbf{Grad-CAM}, ofreciendo finalmente una reflexión sobre las implicaciones éticas y los hallazgos obtenidos en este cruce entre tecnología y criminología.